\subsection{FCC-hh: The energy-frontier collider with the broadest exploration potential}

The physics landscape emerging from the FCC-ee, complemented by the HL-LHC exploration, will significantly raise the bar for the performance requirement of subsequent colliders. 
Several challenges remain open.

\begin{itemize}
\item The extension of Higgs boson property measurements to several processes driven by smaller effective couplings 
(e.g., rare decays such as $\PH \to \PGg\PGg$, $\PGmp\PGmm$, $\PZ\PGg$), 
and to processes requiring higher energy (e.g., $\PQt\PAQt\PH$ and $\PH\PH$ production).
\item The extension of studies of EW dynamics to a regime well beyond the scale of the EW symmetry breaking.
\item The direct exploration of the multi-TeV energy scale, 
to identify and directly study the sources of possible deviations found in the FCC-ee precision measurements, 
and to systematically extend the mass reach for a broad spectrum of BSM theories, 
particularly those whose impact on low-energy EW and Higgs observables is suppressed and those whose primary production processes require initial-state gluons.
\end{itemize}

The 100\,TeV FCC-hh provides the most complete and effective facility to tackle these challenges. 
The study of the FCC-hh physics potential is very mature; it occupied the first phase of the efforts towards the FCC CDR, thanks to a world-wide effort whose results have been documented in extensive reports~\cite{Arkani-Hamed:2015vfh,Mangano:2017tke}, partly summarised in Volumes~1 and~3 of the CDR~\cite{fcc-phys-cdr,FCC-hhCDR}. 
Some of the key results are recalled here. 
Recent preliminary studies of physics performance at different collider energies, in the \mbox{80--120}\,TeV range, have been done in the context of this feasibility study, and are reported in Section~\ref{sec:PhysicsCaseHHscenarios}. 
Section~\ref{sec:PhysicsCaseHHFPF} outlines the opportunities offered by a forward-physics facility, which would expand the landscape of FCC-hh physics measurements presented in the CDR.  


On the Higgs boson side, the data sample produced by FCC-hh of over 20~billion Higgs bosons will bring the absolute determination of the couplings to muons, to photons, to the top quark, and to $\PZ\PGg$ below the per-cent level. 
The large production yield of Higgs bosons at large transverse momentum allows measurements to be performed in kinematic regions with optimal signal-to-background ratios and reduced experimental systematic uncertainties. 
The possibility to accurately measure the ratio of those couplings to couplings precisely measured at the FCC-ee 
(e.g., the $\PH\PZ\PZ^{*}$ and $\PQt\PAQt\PZ$ couplings), 
enables the FCC-hh to deliver sub-percent absolute measurements of those couplings, which are beyond the reach of lepton colliders. 
The large $\PH\PH$ production rate, furthermore, will bring the uncertainty on the Higgs self-coupling below the 5\% level, even with conservative estimates of the experimental systematic uncertainties. 
The studies of Higgs boson production in very-high $Q^2$ processes, whether in direct, associated, or vector boson fusion/scattering production, 
test the existence of higher-dimensional effective operators in ways that are complementary to what is accessible even at the highest-energy lepton colliders. 

The direct search for new particles extends the mass reach to the several tens of TeV range, in particular around 40\,TeV for $s$-channel produced EW or coloured resonances, and in the 10--20\,TeV range for pair-produced strongly-interacting particles, such as squarks, stops, vector-like top partners, and gluinos. 
Weakly interacting particles can be probed up to 1.5--5\,TeV, depending on their weak multiplet assignment, allowing in particular the theoretically limited spectrum of possible WIMP dark-matter candidates to be covered.

Likewise, the search for partners of the Higgs boson extends to the multi-TeV region, a sensitivity that, together with the precise determination of the Higgs boson self-coupling, constrains, in a unique way, extensions of the Higgs sector leading to a strong first-order EW phase transition in the early universe.

Last but not least, the FCC-hh is a unique facility to extend the study of the less known manifestation of the SM dynamics, namely the behaviour of high-density and high-temperature strongly interacting plasmas, 
through the analysis of data collected in central lead-lead collisions at nucleon-nucleon collision energies up to $\sqrt{s} \approx 40$\,TeV~\cite{Dainese:2016gch}. 
Despite the significant progress made over the last decades at the SPS, at RHIC, and at the LHC regarding our understanding of the properties of the hot QCD medium created in high-energy heavy-ion collisions, grounded on increasingly precise and diverse experimental measurements, much remains unknown about the physics behind this most exotic area of QCD.
Details of the concrete physics potential of heavy-ion running at the FCC-hh are provided in Refs.~\cite{Dainese:2016gch,Mangano:2017tke}. 
Aside from its intrinsic and undisputed scientific value, this part of the physics programme, exclusively accessible at hadron colliders, offers the opportunity to enlarge the community of scientists attracted to the FCC, providing them the sole facility that can continue and extend the study of the high-density and high-temperature QCD medium produced in high-energy heavy ion collisions and the associated collective QCD phenomena, presently being actively pursued at the LHC.

The precision programme of the known elements of the SM should be complemented by a direct exploration of the unknown. Exploration requires breadth. 
To understand the properties of the visible Universe, large-scale structure surveys are performed, gravitational waves with pulsar timing and much more are sought. 
In the same vein, to understand the microverse this must be explored with as many different microscopes as possible. 
The FCC-ee will directly explore the evasive and weakly-coupled light new-physics frontier, 
while the FCC-hh will offer a direct way to scan the new energy frontier, above the TeV scale. 
Together, the FCC-ee and the FCC-hh complement each other in this exploration of the unknown in a more comprehensive way than any other energy frontier projects presently under consideration.